\usepackage{tcolorbox}
\usepackage{fontspec}
\usepackage{graphicx}
\usepackage{newunicodechar}

\newtcolorbox{myquote}{left skip=1em, colback=gray!25!white, colframe=gray!25!white}
\renewenvironment{quote}{\begin{myquote}}{\end{myquote}}

% Dedicated fallback font for standalone Unicode math-like glyphs
\newfontfamily\mathfallback{STIX Two Math}

% comparison / operators
\newunicodechar{≈}{\ensuremath{\approx}}
\newunicodechar{∝}{\ensuremath{\propto}}
\newunicodechar{≤}{\ensuremath{\le}}
\newunicodechar{≥}{\ensuremath{\ge}}
\newunicodechar{≠}{\ensuremath{\neq}}
\newunicodechar{±}{\ensuremath{\pm}}
\newunicodechar{×}{\ensuremath{\times}}
\newunicodechar{·}{\ensuremath{\cdot}}

% roots / arrows
% Use a real glyph from a math-capable font; raise slightly for better baseline alignment
\newunicodechar{√}{\raisebox{0.10ex}{{\mathfallback √}}}
\newunicodechar{→}{\ensuremath{\to}}
\newunicodechar{←}{\ensuremath{\leftarrow}}
\newunicodechar{↔}{\ensuremath{\leftrightarrow}}
\newunicodechar{⇒}{\ensuremath{\Rightarrow}}
\newunicodechar{⇔}{\ensuremath{\Leftrightarrow}}

% superscripts often produced by AI/copied text
\newunicodechar{⁰}{\textsuperscript{0}}
\newunicodechar{¹}{\textsuperscript{1}}
\newunicodechar{²}{\textsuperscript{2}}
\newunicodechar{³}{\textsuperscript{3}}
\newunicodechar{⁴}{\textsuperscript{4}}
\newunicodechar{⁵}{\textsuperscript{5}}
\newunicodechar{⁶}{\textsuperscript{6}}
\newunicodechar{⁷}{\textsuperscript{7}}
\newunicodechar{⁸}{\textsuperscript{8}}
\newunicodechar{⁹}{\textsuperscript{9}}
\newunicodechar{⁺}{\textsuperscript{+}}
\newunicodechar{⁻}{\textsuperscript{-}}
\newunicodechar{⁼}{\textsuperscript{=}}
\newunicodechar{⁽}{\textsuperscript{(}}
\newunicodechar{⁾}{\textsuperscript{)}}
